% !TEX program = xelatex
% =============================================================================
% Rapport technique - Projet RESULTATS-BF - Role : Architecte reseau et repartition
% Compilation recommandee : XeLaTeX ou LuaLaTeX (Overleaf : Menu > Compiler > XeLaTeX)
% pdfLaTeX fonctionne aussi (police Libertinus via le paquet "libertinus").
% =============================================================================
\documentclass[11pt,a4paper,oneside]{report}

% ---------------------------------------------------------------- Langue, police
\usepackage{iftex}
\ifPDFTeX
\usepackage[utf8]{inputenc}
\usepackage[T1]{fontenc}
\usepackage{libertinus}
\else
\usepackage{fontspec}
% Polices designees par nom de fichier : fonctionne sur Overleaf, TeX Live et MiKTeX
\setmainfont{LibertinusSerif}[Extension=.otf, UprightFont=*-Regular,
BoldFont=*-Bold, ItalicFont=*-Italic, BoldItalicFont=*-BoldItalic]
\setsansfont{LibertinusSans}[Extension=.otf, UprightFont=*-Regular,
BoldFont=*-Bold, ItalicFont=*-Italic]
% Libertinus Mono n'existe qu'en romain : gras et italique simules
\setmonofont{LibertinusMono}[Extension=.otf, UprightFont=*-Regular,
Scale=0.88, AutoFakeBold=2, AutoFakeSlant=0.15]
\fi
\ifPDFTeX
\usepackage[french]{babel}
\NoAutoSpaceBeforeFDP   % espaces deja saisis a la main ; evite "11 :20"
\else
\usepackage{polyglossia}
\setdefaultlanguage[autospacing=false]{french}  % espaces saisis a la main
\fi

% ---------------------------------------------------------------- Mise en page
\usepackage[a4paper,top=2.5cm,bottom=2.5cm,left=2.5cm,right=2.3cm,headheight=15pt]{geometry}
\usepackage{setspace}\setstretch{1.12}
\usepackage{parskip}
\usepackage{graphicx}
\usepackage[table]{xcolor}
\usepackage{tabularx,booktabs,longtable,array,multirow}
\usepackage{enumitem}\setlist{itemsep=2pt,topsep=4pt}
\usepackage{amssymb}
\usepackage{tikz}
\usetikzlibrary{positioning,arrows.meta,fit,calc,shapes.geometric}
\usepackage{eso-pic}
\usepackage{float}
\usepackage[labelsep=endash,font=small]{caption}
\usepackage[hidelinks]{hyperref}

% ---------------------------------------------------------------- Couleurs (page de garde uniquement)
\definecolor{vert}{RGB}{0,150,57}
\definecolor{bleugarde}{RGB}{0,51,153}
\definecolor{vertclair}{RGB}{200,235,200}

% ---------------------------------------------------------------- En-tete : une ligne, gras italique
\usepackage{fancyhdr}
\pagestyle{fancy}
\fancyhf{}
\fancyhead[L]{\normalfont\textbf{\textit{Projet RÉSULTATS-BF -- Dossier d'architecture}}}
\fancyhead[R]{\normalfont\textbf{\textit{Groupe 7}}}
\fancyfoot[C]{\normalfont\thepage}
\renewcommand{\headrulewidth}{0.4pt}
\fancypagestyle{plain}{\fancyhf{}%
	\fancyhead[L]{\normalfont\textbf{\textit{Projet RÉSULTATS-BF -- Dossier d'architecture}}}%
	\fancyhead[R]{\normalfont\textbf{\textit{Groupe 7}}}%
	\fancyfoot[C]{\normalfont\thepage}\renewcommand{\headrulewidth}{0.4pt}}

% ---------------------------------------------------------------- Titres : Chapitre 1 : / I - / A - / 1 /
\usepackage{titlesec}
\titleformat{\chapter}[hang]{\normalfont\LARGE\bfseries}{Chapitre \thechapter\ :}{0.5em}{}
\titlespacing*{\chapter}{0pt}{-10pt}{18pt}
\renewcommand{\thesection}{\Roman{section}}
\renewcommand{\thesubsection}{\Alph{subsection}}
\renewcommand{\thesubsubsection}{\arabic{subsubsection}}
\titleformat{\section}{\normalfont\Large\bfseries}{\thesection\ -}{0.5em}{}
\titleformat{\subsection}{\normalfont\large\bfseries}{\thesubsection\ -}{0.5em}{}
\titleformat{\subsubsection}{\normalfont\normalsize\bfseries}{\thesubsubsection\ /}{0.5em}{}
\setcounter{secnumdepth}{3}
\setcounter{tocdepth}{2}
\makeatletter
\renewcommand{\l@section}{\@dottedtocline{1}{1.5em}{2.6em}}
\renewcommand{\l@subsection}{\@dottedtocline{2}{4.1em}{2.4em}}
\makeatother

% ---------------------------------------------------------------- Code (noir et blanc)
\usepackage{listings}
\lstset{
	basicstyle=\ttfamily\footnotesize,
	keywordstyle=\bfseries,
	commentstyle=\itshape,
	stringstyle=\ttfamily,
	columns=fullflexible,
	keepspaces=true,
	breaklines=true,
	breakatwhitespace=false,
	postbreak=\mbox{$\hookrightarrow$\space},
	frame=single,
	framerule=0.4pt,
	rulecolor=\color{black},
	numbers=left,
	numberstyle=\tiny,
	numbersep=6pt,
	xleftmargin=1.6em,
	framexleftmargin=1.4em,
	showstringspaces=false,
	tabsize=2,
	extendedchars=true,
	inputencoding=utf8,
	literate={°}{{\textdegree}}1 {é}{{\'e}}1 {è}{{\`e}}1 {à}{{\`a}}1
	{—}{{---}}1 {→}{{$\rightarrow$}}1
}
\lstdefinelanguage{texte}{}
\newcommand{\cmd}[1]{\texttt{#1}}
\setlength{\emergencystretch}{3em}\hbadness=10000

% ---------------------------------------------------------------- Encadres sobres (noir et blanc)
\usepackage[most]{tcolorbox}
\newtcolorbox{aretenir}[1][]{enhanced,breakable,colback=white,colframe=black,
	boxrule=0.6pt,arc=0pt,left=6pt,right=6pt,top=4pt,bottom=4pt,
	title={\textbf{#1}},coltitle=black,colbacktitle=white,
	attach boxed title to top left={xshift=8pt,yshift=-6pt},
	boxed title style={boxrule=0.6pt,arc=0pt}}
\newtcolorbox{soutenance}{enhanced,breakable,colback=white,colframe=black,
	boxrule=0.4pt,arc=0pt,borderline west={2.5pt}{0pt}{black},
	left=8pt,right=6pt,top=4pt,bottom=4pt}


% ---------------------------------------------------------------- Parties a rediger par les autres membres
\newtcolorbox{afairebox}[1]{enhanced,breakable,colback=white,colframe=black,
	boxrule=0.6pt,arc=0pt,borderline={0.4pt}{2pt}{black,dashed},
	left=6pt,right=6pt,top=4pt,bottom=4pt,
	title={\textbf{À rédiger par #1}},coltitle=black,colbacktitle=white,
	attach boxed title to top left={xshift=8pt,yshift=-6pt},
	boxed title style={boxrule=0.6pt,arc=0pt}}
\newcommand{\afaire}[2]{\begin{afairebox}{#1}\small #2\end{afairebox}}
\newcommand{\auteur}[1]{\par\noindent{\small\itshape Rédaction : #1}\par\medskip}
\begin{document}
	\emergencystretch=3em
	\sloppy
	
	% =============================================================================
	%                              PAGE DE GARDE (couleur)
	% =============================================================================
	\begin{titlepage}
		\newgeometry{top=1.4cm,bottom=1.4cm,left=1.6cm,right=1.6cm}
		\thispagestyle{empty}
		\noindent
		\begin{minipage}[c]{0.36\textwidth}\centering\small
			\textbf{Université Joseph KI-ZERBO}\\[4pt]
			\rule{2.2cm}{0.4pt}\\[4pt]
			\textbf{UFR Sciences Exactes et Appliquées}\\[4pt]
			\rule{2.2cm}{0.4pt}\\[4pt]
			\textbf{Laboratoire LAMI}
		\end{minipage}%
		\hfill
		\begin{minipage}[c]{0.24\textwidth}\centering
			\IfFileExists{ujkz.png}{\includegraphics[height=2.8cm]{ujkz.png}}{%
				\fbox{\begin{minipage}[c][2.6cm][c]{2.4cm}\centering\footnotesize LOGO\\(ujkz.png)\end{minipage}}}
		\end{minipage}%
		\hfill
		\begin{minipage}[c]{0.36\textwidth}\centering\small
			\textbf{BURKINA FASO}\\[4pt]
			\rule{2.2cm}{0.4pt}\\[4pt]
			\textit{La Patrie ou la Mort, nous Vaincrons}\\[4pt]
			\rule{2.2cm}{0.4pt}\\[4pt]
			\textbf{Année académique : 2026-2027}
		\end{minipage}
		
		\vspace{0.5cm}
		\noindent{\color{black}\rule{\textwidth}{0.8pt}}\\[-0.68\baselineskip]
		\noindent{\color{bleugarde}\rule{\textwidth}{3pt}}
		
		\vfill
		\begin{center}
			{\fontsize{34}{40}\selectfont\bfseries\color{black}DOSSIER D'ARCHITECTURE}
		\end{center}
		
		\vfill
		\begin{center}
			{\setlength{\fboxrule}{3pt}%
				\fcolorbox{vert}{white}{\parbox{0.92\linewidth}{\centering\vspace{8pt}
						{\Large\bfseries\color{black} Module : Virtualisation et Cloud Computing}\vspace{8pt}}}}
		\end{center}
		
		\vfill
		\begin{center}
			\large
			\begin{tabular}{@{}l l@{}}
				\textbf{Filière :}    & Informatique \\[0.15cm]
				\textbf{Niveau :}     & Master 2 \\[0.15cm]
				\textbf{Spécialité :} & ILSII
			\end{tabular}
		\end{center}
		
		\vfill
		\begin{center}
			\begin{tcolorbox}[enhanced,width=0.92\linewidth,center,arc=0pt,
				colback=white,colframe=vert,boxrule=3pt,
				halign=center,top=16pt,bottom=12pt,
				title={THÈME},fonttitle=\large\bfseries,coltitle=black,
				attach boxed title to top center={yshift=-\tcboxedtitleheight/2},
				boxed title style={colback=white,colframe=vert,boxrule=2pt,arc=0pt}]
				{\Large\bfseries Projet RÉSULTATS-BF : infrastructure résiliente et scalable,
					entièrement décrite en code avec Vagrant}
			\end{tcolorbox}
		\end{center}
		
		\vfill
		\noindent
		\begin{minipage}[t]{0.55\textwidth}
			\textbf{\large Présenté par :}\\[6pt]
			\begin{tabular}{@{}l@{}}
				BANCE François Xavier \\[3pt]
				DAMA Soumana \\[3pt]
				PANANDETIGUI Hawa \\[3pt]
				DAMBRE Brigitte
			\end{tabular}
		\end{minipage}%
		\hfill
		\begin{minipage}[t]{0.40\textwidth}\raggedleft
			\textbf{\large Encadrant :}\\[6pt]
			\large Dr Flavien Hervé SOMDA
		\end{minipage}
		
		\vfill
		\begin{center}{\large Novembre 2026}\end{center}
	\end{titlepage}
	\restoregeometry
	
	\pagenumbering{roman}
	\tableofcontents
	\clearpage
	\pagenumbering{arabic}
	
	% =============================================================================
	\chapter{Contexte et objectifs}
	\auteur{Le groupe}
	
	\section{Le problème}
	
	Chaque année, le site de consultation des résultats d'examen tombe dans les minutes qui suivent
	leur publication. Le programme n'est pas en cause : c'est son infrastructure, qui n'a été pensée
	ni pour absorber un pic de visites, ni pour survivre à la panne d'une seule machine.
	
	\section{Les objectifs}
	
	La plateforme doit répondre à deux exigences :
	\begin{itemize}
		\item \textbf{résilience} : la panne d'une machine n'interrompt pas le service ;
		\item \textbf{scalabilité} : on absorbe plus de trafic en ajoutant des serveurs, sans réécrire
		la configuration.
	\end{itemize}
	Toute l'infrastructure est \textbf{décrite en code} : la commande \cmd{vagrant destroy -f}
	suivie de \cmd{vagrant up} reconstruit la plateforme entière, sans aucun geste manuel.
	
	\section{Organisation du groupe}
	
	\begin{table}[H]\centering\small
		\begin{tabularx}{\textwidth}{@{}l l X@{}}
			\toprule
			\textbf{Rôle} & \textbf{Membre} & \textbf{Périmètre} \\
			\midrule
			Architecte réseau et répartition & BANCE François Xavier & Plan d'adressage, Vagrantfile, répartiteur, sonde de santé, extension E1 \\
			Ingénieur applicatif & DAMA Soumana & Service HTTP, absence d'état, gestion des erreurs de base \\
			Ingénieur données & PANANDETIGUI Hawa & PostgreSQL, jeu de données, réplication, procédure de bascule \\
			Ingénieur essais et observabilité & DAMBRE Brigitte & Scénarios V1 à V6, campagne de mesures \\
			\bottomrule
		\end{tabularx}
	\end{table}
	
	% =============================================================================
	\chapter{Architecture}
	\auteur{BANCE François Xavier}
	
	\section{Vue d'ensemble}
	
	\begin{figure}[H]\centering
		\begin{tikzpicture}[
			noeud/.style={draw,thick,minimum width=2.3cm,minimum height=1.05cm,align=center,font=\small},
			fleche/.style={-{Latex[length=2.5mm]},thick}]
			\node[noeud,dashed,minimum width=4.2cm] (hote) at (0,0) {Visiteurs (hôte)};
			\node[draw,thick,ellipse,minimum width=3.6cm,font=\small,align=center] (vip) at (0,-1.8) {IP virtuelle\\192.168.62.10};
			\node[noeud] (lb1) at (-3,-3.7) {\textbf{lb1} (.11)\\HAProxy\\maître};
			\node[noeud] (lb2) at (3,-3.7) {\textbf{lb2} (.12)\\HAProxy\\secours};
			\node[noeud] (app1) at (-4,-6.1) {\textbf{app1}\\.21:8000};
			\node[noeud] (app2) at (0,-6.1) {\textbf{app2}\\.22:8000};
			\node[noeud] (appn) at (4,-6.1) {\textbf{appN}\\.2N:8000};
			\node[noeud] (bdd1) at (-2.6,-8.7) {\textbf{bdd1} (.31)\\PostgreSQL primaire};
			\node[noeud] (bdd2) at (3.4,-8.7) {\textbf{bdd2} (.32)\\réplique};
			\draw[fleche] (hote) -- (vip);
			\draw[fleche] (vip) -- node[above left,font=\footnotesize]{porte la VIP} (lb1);
			\draw[fleche,dashed] (vip) -- node[above right,font=\footnotesize]{si lb1 tombe} (lb2);
			\draw[<->,dotted,thick] (lb1) -- node[above,font=\footnotesize]{VRRP, 1 s} (lb2);
			\foreach \a in {app1,app2,appn}{\draw[fleche] (lb1.south) -- (\a.north);}
			\foreach \a in {app1,app2,appn}{\draw[fleche] (\a.south) -- (bdd1.north);}
			\draw[fleche] (bdd1) -- node[below,font=\footnotesize]{réplication en flux} (bdd2);
		\end{tikzpicture}
		\caption{Architecture de la plateforme (extension E1 activée)}
	\end{figure}
	
	Les visiteurs n'utilisent qu'une adresse, l'\textbf{IP virtuelle}, portée par le répartiteur maître.
	Celui-ci distribue les requêtes entre les serveurs applicatifs, qui interrogent la base primaire.
	La réplique suit la primaire en continu et peut être promue en cas de panne.
	
	\section{Plan d'adressage}
	
	Le groupe 7 utilise le sous-réseau imposé \cmd{192.168.(55+7).0/24} = \textbf{\cmd{192.168.62.0/24}},
	dans la plage autorisée par VirtualBox sans modification du poste.
	
	\begin{table}[H]\centering\small
		\caption{Plan d'adressage}
		\begin{tabular}{@{}l l l l@{}}
			\toprule
			\textbf{Machine} & \textbf{Rôle} & \textbf{Adresse} & \textbf{RAM} \\
			\midrule
			(VIP) & Adresse du service, partagée lb1/lb2 & 192.168.62.10 & -- \\
			lb1 / lb2 & Répartiteurs HAProxy (lb2 : extension E1) & .11 / .12 & 512 Mo \\
			app1 \ldots\ appN & Serveurs applicatifs, port 8000 ($2 \le N \le 9$) & .21 \ldots\ .2N & 512 Mo \\
			bdd1 / bdd2 & PostgreSQL primaire / réplique & .31 / .32 & 768 Mo \\
			\bottomrule
		\end{tabular}
	\end{table}
	
	Chaque famille de machines occupe sa dizaine (\cmd{.1x}, \cmd{.2x}, \cmd{.3x}) ; l'adresse du
	serveur $i$ vaut $20+i$. Seul lb1 expose un port vers l'hôte (\cmd{localhost:8080}), lié à
	\cmd{127.0.0.1} : tout le reste du trafic circule sur le réseau privé (R8). Le socle (5 VM)
	occupe 3 Go de mémoire, 3,5 Go avec l'extension E1.
	
	\section{La source unique de vérité}
	
	La liste des machines (nom, adresse, mémoire, rôle) est écrite \textbf{une seule fois}, dans une
	table du Vagrantfile. Cette table sert à créer les VM \emph{et} est écrite dans
	\cmd{config/noeuds.yml}, que les scripts lisent \emph{dans} les VM pour \textbf{engendrer} les
	configurations de HAProxy et de keepalived à partir de modèles (\cmd{config/*.tpl}).
	Aucune adresse n'est écrite en dur : pour passer de 2 à 4 serveurs, on ne change qu'une valeur,
	\cmd{NB\_APP} (R3).
	
	Le Vagrantfile lit ses paramètres et ses secrets dans \cmd{.env}, jamais versionné ; le modèle
	\cmd{.env.exemple}, versionné, ne contient aucun mot de passe (R7). Il refuse de démarrer si une
	valeur est incohérente (groupe hors de 1 à 8, $N$ hors de 2 à 9, mot de passe absent).
	
	% =============================================================================
	\chapter{Justification des choix}
	
	\section{Infrastructure et répartition}
	\auteur{BANCE François Xavier}
	
	\begin{longtable}{@{}p{0.27\textwidth} p{0.67\textwidth}@{}}
		\toprule
		\textbf{Choix} & \textbf{Justification} \\
		\midrule
		\endhead
		Box \cmd{bento/ubuntu-24.04} épinglée (\cmd{202510.26.0}) & Même système pour tous les membres : sans épinglage, un coéquipier pourrait recevoir une autre version et la plateforme ne serait plus reproductible. \\
		\addlinespace
		Clones liés, 1 vCPU par VM & Création rapide et économe en disque sur un poste de 8 Go avec un processeur 2 cœurs. \\
		\addlinespace
		\textbf{HAProxy} comme répartiteur & Sonde de santé HTTP réglable à la seconde, répartition en tourniquet, vérification d'une configuration avant application (\cmd{haproxy -c}), rechargement sans couper les connexions. \\
		\addlinespace
		Sonde \cmd{GET /sante}, \cmd{inter 1s fall 2 rise 2} & Un serveur est écarté après deux échecs consécutifs, soit environ 2 s, sous les 5 s exigées (R4) ; il est réintégré automatiquement après deux succès. La sonde vérifie la réponse HTTP (couche 7), pas seulement le port. \\
		\addlinespace
		\cmd{timeout check 3s} & Quand bdd1 est en panne, \cmd{/sante} attend 2 s sur bdd1 (minimum imposé par libpq) avant d'essayer bdd2. Avec 1 s, tous les serveurs seraient écartés à tort après la bascule de la base. La connexion de la sonde reste limitée à 1 s : un serveur éteint est toujours écarté en environ 2 s. \\
		\addlinespace
		\cmd{retries 2} et \cmd{option redispatch} & Un échec de \emph{connexion} est rejoué sur un autre serveur : moins de requêtes perdues lors d'une panne (mesuré en V2). \\
		\addlinespace
		\textbf{keepalived (VRRP)} pour E1 & Une IP virtuelle partagée supprime le point de défaillance unique du répartiteur sans rien changer côté visiteurs. Annonces en \emph{unicast} (plus fiable que la multidiffusion sur un réseau VirtualBox), sonde locale de HAProxy (un répartiteur dont HAProxy est arrêté cède l'adresse). \\
		\addlinespace
		\cmd{advert\_int 1} & Compromis entre rapidité de bascule (environ 3,6 s en panne brutale) et risque de fausse bascule sur un poste chargé. \\
		\addlinespace
		E1 activable (\cmd{AVEC\_LB2}) & Le socle reste à 5 VM par défaut ; l'extension n'est allumée que lorsqu'elle est utile. \\
		\addlinespace
		Scripts idempotents & Chaque script compare le fichier engendré à l'existant (\cmd{cmp}) et ne recharge un service que s'il a changé ; les blocs ajoutés à des fichiers système sont délimités par des marqueurs. Le répartiteur peut ainsi être reprovisionné à chaque \cmd{vagrant up} (\cmd{run: "always"}) et suivre \cmd{NB\_APP}. \\
		\bottomrule
	\end{longtable}
	
	\section{Application}
	\auteur{DAMA Soumana}
	
	\begin{longtable}{@{}p{0.27\textwidth} p{0.67\textwidth}@{}}
		\toprule
		\textbf{Choix} & \textbf{Justification} \\
		\midrule
		\endhead
		Flask servi par \textbf{gunicorn}, 3 processus & Flask suffit pour trois routes JSON ; gunicorn sert les requêtes en parallèle. Trois processus sur une VM à 1 vCPU : un compromis entre parallélisme et mémoire (512 Mo). Service \cmd{systemd} sous l'utilisateur restreint \cmd{www-data}, redémarrage automatique (\cmd{Restart=always}). \\
		\addlinespace
		\textbf{Aucun état local} & Ni session, ni fichier, ni compteur en mémoire : toute donnée est dans PostgreSQL. N'importe quel serveur peut traiter n'importe quelle requête, ce qui rend l'ajout de serveurs immédiat (R3). \\
		\addlinespace
		\cmd{/sante} exécute \cmd{SELECT 1} & La sonde vérifie le service \emph{et} son accès à la base : un serveur qui ne joint plus la base est écarté par HAProxy. \\
		\addlinespace
		Erreurs 400 / 404 / 503 en JSON & Paramètre manquant (400), matricule inconnu (404), base indisponible (503). L'erreur brute de PostgreSQL est journalisée côté serveur, jamais renvoyée au visiteur (V4). Requêtes SQL paramétrées : pas d'injection SQL. \\
		\addlinespace
		\cmd{host=bdd1,bdd2} et \cmd{target\_session\_attrs} = \cmd{read-write} & libpq choisit automatiquement la base qui accepte les écritures, donc la primaire, y compris après une bascule. \\
		\addlinespace
		Mémorisation de la dernière primaire & Chaque processus essaie d'abord la dernière base qui a répondu. Sans cela, chaque requête attendait 2 s sur bdd1 éteinte après une bascule (défaut révélé par V4, corrigé). \\
		\addlinespace
		Secrets dans \cmd{app.env} & Fichier \cmd{/etc/resultats-bf/app.env}, engendré par \cmd{app.sh} à partir des variables transmises par Vagrant, droits 640 (\cmd{root:www-data}) ; aucun secret dans le code ni dans Git (R7). \\
		\bottomrule
	\end{longtable}
	
	\section{Données}
	\auteur{PANANDETIGUI Hawa}
	
	\begin{longtable}{@{}p{0.27\textwidth} p{0.67\textwidth}@{}}
		\toprule
		\textbf{Choix} & \textbf{Justification} \\
		\midrule
		\endhead
		PostgreSQL 16 (paquet Ubuntu) & Version fournie par Ubuntu 24.04 : installation reproductible, sans dépôt externe. \\
		\addlinespace
		Jeu de données déterministe & 50~000 candidats (\cmd{BF00001} à \cmd{BF50000}) engendrés par \cmd{bdd.sh} : la même base est recréée à chaque reconstruction. \\
		\addlinespace
		\textbf{Réplication en flux} & bdd2 est initialisée par \cmd{pg\_basebackup} et suit bdd1 en continu (\cmd{wal\_level = replica}, rôle \cmd{replicateur} dédié). Elle est en lecture seule (\cmd{pg\_is\_in\_recovery() = true}) jusqu'à sa promotion. \\
		\addlinespace
		Données hors de \cmd{/vagrant} & Les fichiers de PostgreSQL restent dans la VM : le dossier partagé avec l'hôte ne sert qu'aux scripts. \\
		\addlinespace
		Accès limités (\cmd{pg\_hba.conf}) & Seul le réseau privé \cmd{192.168.62.0/24} est autorisé, avec authentification \cmd{scram-sha-256}. bdd2 reçoit dès sa création l'accès de l'application, pour être utilisable dès sa promotion. \\
		\addlinespace
		\textbf{Promotion manuelle} (\cmd{pg\_promote()}) & Une promotion automatique sans arbitre risquerait deux primaires (\emph{split-brain}). La procédure manuelle est écrite et chronométrée (R6). \\
		\bottomrule
	\end{longtable}
	
	% =============================================================================
	\chapter{Procédures d'exploitation}
	
	\section{Reconstruire la plateforme}
	\auteur{BANCE François Xavier}
	
	\begin{lstlisting}[language=texte,numbers=none]
		copy .env.exemple .env        # puis remplacer les deux mots de passe
		vagrant up                    # 13 min 22 s mesurees sur notre poste de reference
	\end{lstlisting}
	
	Après une destruction (\cmd{vagrant destroy -f}), vérifier qu'aucun dossier
	\cmd{resultats-bf-g7-*} ne subsiste dans \cmd{VirtualBox VMs} : sous Windows, VirtualBox peut en
	laisser, et ils empêchent la recréation de la VM (incident rencontré et résolu, voir les limites).
	
	\section{Ajouter ou retirer un serveur applicatif}
	\auteur{BANCE François Xavier}
	
	\textbf{Ajouter} (exemple : passer à 4 serveurs) :
	\begin{lstlisting}[language=texte,numbers=none]
		$env:NB_APP=4; vagrant up; Remove-Item Env:NB_APP     # PowerShell
		NB_APP=4 vagrant up                                   # bash
	\end{lstlisting}
	\cmd{vagrant up} crée app3 et app4, puis rejoue \cmd{lb.sh} sur le répartiteur : la configuration
	de HAProxy est régénérée avec les quatre serveurs. Pour un changement durable, écrire
	\cmd{NB\_APP=4} dans \cmd{.env}.
	
	\textbf{Retirer} : toujours \emph{retirer du répartiteur d'abord}, détruire ensuite. Dans l'ordre
	inverse, HAProxy enverrait des requêtes à une machine détruite pendant environ 2 s.
	\begin{lstlisting}[language=texte,numbers=none]
		vagrant provision lb1                                          # 1. lb1 revient a 2 serveurs
		$env:NB_APP=4; vagrant destroy -f app3 app4; Remove-Item Env:NB_APP   # 2. destruction
	\end{lstlisting}
	
	\section{Bascule du répartiteur (extension E1)}
	\auteur{BANCE François Xavier}
	
	La bascule est \textbf{automatique} : aucune commande n'est nécessaire. Pour l'observer :
	\begin{lstlisting}[language=texte,numbers=none]
		curl.exe http://192.168.62.10/repartiteur      # repond lb1, ou lb2 apres une panne de lb1
	\end{lstlisting}
	Au retour de lb1, sa priorité plus forte (150 contre 100) lui fait reprendre l'IP virtuelle
	(préemption). Le rôle de chaque machine se lit dans le journal :
	\cmd{journalctl -u keepalived} (\emph{Entering MASTER STATE} / \emph{BACKUP STATE}).
	
	\section{Bascule de la base de données (R6)}
	\auteur{PANANDETIGUI Hawa}
	
	\begin{enumerate}
		\item \textbf{Constater la panne} : \cmd{curl.exe -i http://localhost:8080/sante} répond \cmd{503}.
		\item \textbf{Promouvoir bdd2} :
		\begin{lstlisting}[language=texte,numbers=none]
			vagrant ssh bdd2 -c "sudo -u postgres psql -c 'SELECT pg_promote();'"
		\end{lstlisting}
		\item \textbf{Vérifier} : \cmd{SELECT pg\_is\_in\_recovery();} sur bdd2 renvoie \cmd{f} (primaire) ;
		\cmd{/sante} répond de nouveau \cmd{200}. L'application se reconnecte seule, sans modification.
		\item \textbf{Revenir à l'état normal} : ne \emph{jamais} redémarrer bdd1 telle quelle (elle
		repartirait primaire avec des données périmées : deux primaires). On reconstruit les deux VM :
		\cmd{vagrant destroy -f bdd1 bdd2} puis \cmd{vagrant up bdd1 bdd2}.
	\end{enumerate}
	\textbf{Durée mesurée (V4)} : 22,5 s d'interruption, de la panne de bdd1 à la reprise du service.
	
	\section{Déploiement sans coupure (extension E4)}
	\auteur{BANCE François Xavier et DAMA Soumana}
	
	Après modification de \cmd{app/app.py}, une seule commande met à jour tous les serveurs :
	\begin{lstlisting}[language=texte,numbers=none]
		powershell -ExecutionPolicy Bypass -File scripts\deployer.ps1
	\end{lstlisting}
	Pour chaque serveur, l'un après l'autre, le script : (1) le retire de HAProxy
	(\cmd{set server serveurs\_app/appN state maint} sur le socket d'administration, sur chaque
	répartiteur allumé) ; (2) attend 3 s la fin des requêtes en cours ; (3) le met à jour
	(\cmd{vagrant provision appN --provision-with app}) ; (4) attend que \cmd{/sante} réponde 200 ;
	(5) le remet en service (\cmd{state ready}). La liste des serveurs est lue dans
	\cmd{config/noeuds.yml} : le script suit \cmd{NB\_APP}. Il refuse de retirer un serveur si aucun
	autre n'est sain, et s'arrête en laissant hors service un serveur qui ne redémarre pas.
	
	% =============================================================================
	\chapter{Résultats des scénarios V1 à V6}
	
	Chaque mesure a été faite avec l'outil \cmd{tests/surveiller.ps1}, qui interroge le service
	environ quatre fois par seconde, horodate chaque réponse (une réponse HTTP d'erreur compte comme
	un échec) et calcule requêtes perdues, disponibilité et durée des coupures (précision d'environ
	$\pm 0{,}5$ s). Les sorties brutes sont dans le dossier \cmd{mesures/}.
	
	\section{Synthèse}
	
	\begin{tabularx}{\textwidth}{@{}l X l l@{}}
		\toprule
		\textbf{N°} & \textbf{Mesure relevée} & \textbf{Résultat} & \textbf{Exigence} \\
		\midrule
		V1 & Écart entre serveurs (20 appels) & app1 10, app2 10 : écart 0 \% & $<$ 20 \% \quad\checkmark \\
		V2 & Requêtes perdues ; période dégradée & 1 perdue sur 328 ; coupure 1,16 s & $<$ 5 s \quad\checkmark \\
		V3 & Requêtes par seconde (\cmd{ab -n 2000 -c 50}) & $N=2$ : 26,1 ; $N=4$ : 23,4 & voir V3 \quad $\times$ \\
		V4 & 503 propre ; durée de la bascule & 503 ; 22,5 s & mesurée \quad\checkmark \\
		V5 & Reconstruction complète & 11 min 51 s, code 0 & sans intervention \quad\checkmark \\
		V6 & Reprise après bascule VRRP & 0,80 s ; retour de lb1 : 0,76 s & même adresse \quad\checkmark \\
		E4 & Déploiement pendant la mesure & 583 requêtes, 0 échec ; 121 s & 0 erreur \quad \checkmark \\
		R5 & Écriture sur bdd1, lecture sur bdd2 & ligne lue ; bdd2 refuse l'écriture & réplication \quad\checkmark \\
		\bottomrule
	\end{tabularx}
	
	
	La campagne officielle a été menée par Brigitte sur son propre PC : elle prouve en même temps
	que la plateforme se reconstruit sur le poste d'un autre membre. Les tests d'intégration faits
	sur le poste de François (8 Go) donnent des valeurs concordantes (V1 50/50, V2 environ 2 s,
	V5 13 min 22 s, V6 1,00 s et 3,46 s en coupure brutale) : les résultats sont reproductibles.
	
	\section{V1  Répartition}
	\auteur{DAMBRE Brigitte}
	
	\textbf{Commande :} 20 appels successifs à \cmd{http://localhost:8080/info}
	(\cmd{mesures/V1-20-appels.txt}).\\
	\textbf{Résultat :} app1 = 10, app2 = 10, en alternance stricte ; écart 0 \%.\\
	\textbf{Commentaire :} comportement attendu de \cmd{balance roundrobin} avec deux serveurs de même
	poids.
	
	\section{V2  Panne d'un serveur applicatif}
	\auteur{DAMBRE Brigitte}
	
	\textbf{Commande :} \cmd{tests/surveiller.ps1} sur \cmd{/info} pendant 150 s ;
	\cmd{vagrant halt app1} après 10 s, puis \cmd{vagrant up app1} (\cmd{mesures/V2-panne-app1.txt},
	\cmd{mesures/V2-journal-haproxy.txt}).\\
	\textbf{Résultat :} 328 requêtes, 1 perdue (99,70 \%), une seule coupure de 1,16 s. Journal :
	\emph{app1 is DOWN} (\emph{Layer4 connection refused}) à 23:13:22, puis \emph{app1 is UP}
	(\emph{Layer7 check passed, code 200}) à 23:14:38, sans intervention.\\
	\textbf{Commentaire :} aucune erreur au-delà de 5 s. La requête perdue était déjà en cours sur
	app1 au moment de l'arrêt : \cmd{redispatch} ne rejoue que les échecs de connexion.
	
	\section{V3  Montée en charge}
	\auteur{DAMBRE Brigitte}
	
	\textbf{Commande :} \cmd{ab -n 2000 -c 50}, lancé depuis bdd2 vers lb1, avec $N=2$ puis $N=4$
	(\cmd{mesures/V3-ab-*.txt}).
	
	\begin{table}[H]\centering\small
		\begin{tabular}{@{}l r r r r@{}}
			\toprule
			\textbf{Page} & \multicolumn{2}{c}{\textbf{Requêtes / s}} & \multicolumn{2}{c}{\textbf{Temps moyen (ms)}} \\
			& $N=2$ & $N=4$ & $N=2$ & $N=4$ \\
			\midrule
			\cmd{/resultat} (avec la base) & 26,09 & 23,36 & 1916 & 2140 \\
			\cmd{/info} (sans la base) & 117,21 & 100,14 & 427 & 499 \\
			\bottomrule
		\end{tabular}
	\end{table}
	Aucun échec dans les quatre mesures.\\
	\textbf{Commentaire honnête :} le débit \emph{baisse} avec 4 serveurs. Le test sur \cmd{/info},
	qui n'interroge pas la base, baisse lui aussi : le goulot est le \textbf{processeur du PC},
	partagé par toutes les VM. Ajouter des VM sur la même machine ajoute de la concurrence, pas de
	puissance. Par ailleurs, \cmd{/resultat} est 4,5 fois plus lent que \cmd{/info} : l'application
	ouvre une connexion PostgreSQL à chaque requête. Sur des machines distinctes, la répartition
	augmenterait le débit ; un \emph{pool} de connexions l'améliorerait aussi.
	
	\section{V4  Panne de la base primaire}
	\auteur{DAMBRE Brigitte et PANANDETIGUI Hawa}
	
	\textbf{Commande :} \cmd{tests/surveiller.ps1} sur \cmd{/resultat?matricule=BF00001} pendant
	300 s ; \cmd{vagrant halt bdd1}, vérification de \cmd{/sante}, puis promotion de bdd2
	(\cmd{mesures/V4-apres-correction.txt}).\\
	\textbf{Résultat :} pendant la panne, \cmd{/sante} répond \cmd{503} proprement. Après la promotion,
	le service reprend : interruption totale de \textbf{22,5 s} (procédure manuelle comprise),
	disponibilité de 91,67 \% sur les 5 minutes.\\
	\textbf{Historique :} les premiers essais (fichiers conservés) ont révélé un défaut : après la
	bascule, chaque requête attendait 2 s, car l'application essayait toujours bdd1 en premier.
	Soumana l'a corrigé (mémorisation de la dernière primaire) ; le dernier essai valide la
	correction. Deux micro-coupures (0,8 s et 1,5 s) apparaissent environ 90 s après la reprise ;
	hypothèse : l'arrêt final de la VM bdd1 par Vagrant.
	
	\section{V5  Reconstruction}
	\auteur{DAMBRE Brigitte}
	
	\textbf{Commande :} \cmd{vagrant up} chronométré, sur un clone neuf du dépôt
	(\cmd{mesures/V5-reconstruction-PC-Brigitte.txt}).\\
	\textbf{Résultat :} \textbf{11 min 51 s}, code de retour 0, aucune intervention, sur le PC de
	Brigitte (13 min 22 s sur celui de François). Un \cmd{vagrant provision} rejoué ne modifie rien
	(R2).\\
	\textbf{Commentaire :} la box était déjà présente sur le poste ; sur une machine vierge, ajouter
	son téléchargement (environ 1 Go). Le temps est dominé par le démarrage séquentiel des VM et les
	téléchargements de paquets.
	
	\section{V6  Panne du répartiteur (extension E1)}
	\auteur{BANCE François Xavier et DAMBRE Brigitte}
	
	\begin{table}[H]\centering\small
		\begin{tabularx}{\textwidth}{@{}l l l X@{}}
			\toprule
			\textbf{Situation} & \textbf{Coupure} & \textbf{Requêtes perdues} & \textbf{Explication} \\
			\midrule
			\cmd{vagrant halt lb1} & 0,80 s & 1 / 252 & keepalived annonce son départ (priorité 0) : lb2 devient maître aussitôt \\
			Coupure brutale (\cmd{poweroff}) & 3,46 s & 6 / 162 & lb2 constate l'absence d'annonces (poste de François) \\
			Retour de lb1 & 0,76 s & 1 / 637 & préemption sans perte, puis brève coupure due à Vagrant (voir limites) \\
			\bottomrule
		\end{tabularx}
	\end{table}
	
	En panne brutale, VRRP déclare le maître absent après le \emph{Master Down Interval} :
	$3 \times \text{advert\_int} + (256 - 100)/256 \approx 3{,}6$ s. La mesure (3,46 s) concorde avec
	la théorie. Le service est resté joignable à la même adresse dans tous les cas.
	
	\section{E4  Déploiement sans coupure}
	\auteur{BANCE François Xavier}
	
	\textbf{Commande :} \cmd{tests/surveiller.ps1} sur \cmd{/info} pendant le déploiement de la
	version 1.2 par \cmd{scripts/deployer.ps1}, lancé pendant la mesure (\cmd{mesures/E4-deploiement.txt}).\\
	\textbf{Résultat :} déploiement de la version 1.1 à la version 1.2 en 121 s ; \textbf{583 requêtes,
		0 échec, 100 \% de disponibilité}. Les réponses passent en version 1.2 serveur par serveur :
	app1 à 20:56:36, app2 à 20:57:21. Le déploiement installait aussi l'arrêt propre de gunicorn
	(\cmd{--graceful-timeout 10}).\\
	\textbf{Commentaire :} contrairement à V2, le serveur est retiré du répartiteur \emph{avant}
	d'être arrêté : HAProxy ne lui envoie plus de requête, aucune n'est perdue. Un essai manuel
	préalable (retrait et remise d'app1) donnait déjà 344 requêtes et 0 échec
	(\cmd{mesures/E4-essai-manuel.txt}). Lors d'un essai antérieur, une coupure isolée de 4,6 s
	a été observée sur les deux serveurs à la fois (\cmd{mesures/E4-deploiement-v1.1.txt}) ; sa cause
	n'est pas établie (hypothèse : saturation du processeur de l'hôte) et elle ne s'est pas reproduite.
	
	% =============================================================================
	\chapter{Points de défaillance unique}
	\auteur{Le groupe}
	
	\begin{longtable}{@{}p{0.2\textwidth} p{0.36\textwidth} p{0.36\textwidth}@{}}
		\toprule
		\textbf{Élément} & \textbf{Effet de sa panne} & \textbf{Parade mise en place} \\
		\midrule
		\endhead
		Serveur applicatif & Requêtes en cours perdues & Sonde de santé : éviction en $\approx$ 2 s, réintégration automatique (V2) \\
		\addlinespace
		Répartiteur & Plus aucun accès au site & Extension E1 : lb2 reprend l'IP virtuelle en 1 à 3,5 s (V6) \\
		\addlinespace
		Base primaire bdd1 & \cmd{/sante} en 503, service indisponible & Réplique bdd2 promue manuellement : reprise en 22,5 s (R6, V4) \\
		\addlinespace
		Poste hôte & Toute la plateforme s'arrête & Aucune : limite assumée d'un environnement de laboratoire \\
		\bottomrule
	\end{longtable}
	
	% =============================================================================
	\chapter{Limites et ce que nous ferions autrement}
	
	\section{Infrastructure et répartition}
	\auteur{BANCE François Xavier}
	
	\begin{itemize}
		\item \textbf{Requête en cours perdue (V2)} : une requête déjà transmise à un serveur qui
		tombe est perdue. Pour une maintenance planifiée, c'est évité : \cmd{deployer.ps1} retire le
		serveur du répartiteur avant de l'arrêter (E4, 0 erreur).
		\item \textbf{Petite coupure au retour de lb1} (0,79 s, 10 s après la préemption) : à ce moment,
		Vagrant reconfigure la carte réseau \cmd{eth1} et keepalived quitte brièvement l'état maître
		(journaux de \cmd{systemd-networkd} et de keepalived). Effet propre à Vagrant ; l'option
		keepalived \cmd{preempt\_delay} ferait attendre lb1 avant de reprendre le rôle de maître.
		\item \textbf{Bascule plus lente en panne brutale} (3,46 s contre 1,00 s) : inhérente à VRRP ; un
		\cmd{advert\_int} plus faible l'accélérerait au prix d'un risque de fausse bascule.
		\item \textbf{Ports SSH} : Vagrant ouvre un port SSH par VM, lié à \cmd{127.0.0.1} ; c'est son
		canal d'administration, aucun port applicatif n'est exposé hors de lb1.
		\item \textbf{Dossiers orphelins de VirtualBox sous Windows} : après un \cmd{destroy}, un dossier
		de VM peut subsister et bloquer la recréation ; il faut le vérifier et le supprimer.
		\item \textbf{Arrêt propre des VM lent} : sur notre poste, Vagrant finit souvent par forcer l'arrêt
		après 60 s ; PostgreSQL s'en remet grâce à son journal, mais la cause reste à analyser.
		\item \textbf{Téléchargements répétés} : chaque VM télécharge ses propres paquets ; un cache de
		paquets partagé réduirait la durée de reconstruction, surtout avec une connexion lente.
		\item \textbf{Horloges} : les horloges des VM retardent de 1 à 2 s sur l'hôte ; les durées sont
		donc toujours mesurées côté visiteur, sur une seule horloge.
	\end{itemize}
	
	\section{Application, données et mesures}
	\auteur{DAMA Soumana, PANANDETIGUI Hawa et DAMBRE Brigitte}
	
	\begin{itemize}
		\item \textbf{Une connexion à la base par requête} : coûteuse (V3 : \cmd{/resultat} 4,5 fois plus
		lent que \cmd{/info}). Un \emph{pool} de connexions par processus améliorerait le débit.
		\item \textbf{Bascule de la base manuelle} : 22,5 s d'interruption, dépendant de la réactivité de
		l'opérateur. Un gestionnaire de bascule avec arbitre (par exemple Patroni) l'automatiserait sans
		risque de deux primaires.
		\item \textbf{Retour à la normale par reconstruction} : après une bascule, bdd1 n'est pas
		réintégrée comme réplique ; on reconstruit les deux bases. \cmd{pg\_rewind} permettrait de la
		réintégrer, mais les écritures faites entre-temps doivent être tranchées explicitement.
		\item \textbf{Débit limité par le poste unique} (V3) : toutes les VM partagent le même processeur ;
		la scalabilité horizontale ne se mesure honnêtement que sur des machines distinctes.
		\item \textbf{Démarrage parfois trop lent d'une VM} : lorsque plusieurs VM démarrent, l'une d'elles
		peut dépasser le délai d'attente de Vagrant ; on la relance seule (\cmd{vagrant up app2}).
	\end{itemize}
	
\end{document}